\documentclass{article}
\usepackage[utf8]{inputenc}

\usepackage[lmargin=1in, rmargin=1in, top=0.5in, bottom=0.5in]{geometry}
\usepackage{graphicx}

\title{\textbf{Drew University - CSCI 260 - Fall 2022\\HW6 - C Programming - String and File I/O}}
\date{}

\begin{document}
\maketitle

\vspace{-0.5in}
\textbf{HW Info.}: Due on 12/09/2022 (Fri.) by 11:59pm, 20 points

\textbf{Submission}: Submit a \textit{.zip} file that contains your answer to Problem 1 in \textit{.pdf} and your C problem \textit{Palindrome.c} to Problem 2.

\textbf{Note}: Please follow Drew University's Standard of Academic Integrity and refer to the course syllabus for late submission policy.\\



\begin{enumerate}
    \item (10 points) For the program below, find the value of \textit{int result = count(20);} by either drawing the flow chart similar to that in Figure 17.2 from the textbook (also on Page 7 from slide deck) or the stack operations similar to that Figure 17.14 from the textbook (also on Page 11-12 from slide deck).
    
    \hspace{0.5in}\textit{int count(int arg) \{}
    
    \hspace{0.5in}\textit{\{}
    
    \hspace{0.65in}\textit{if (arg $<$ 1)}
    
    \hspace{0.8in}\textit{return 0;}
    
    \hspace{0.65in}\textit{else if (arg $\%$ 2 == 1)}
    
    \hspace{0.8in}\textit{return (1 + count(arg - 2));}

    \hspace{0.65in}\textit{else}
    
    \hspace{0.8in}\textit{return (1 + count(arg - 1));}
    
    \hspace{0.5in}\textit{\}}
    
    
    \vspace{0.3in}
    \item (10 points) In C, write a \textbf{recursive} function \textit{int Palindrome(char mystr[], int len)} that checks whether the string \textit{char mystr[] = "... ..."} is a \textbf{palindrome} or not. If it is then print ``Palindrome" and if not print ``Not palindrome" to the console.
    
    A string being a palindrome has 2 cases: $(i)$ if the string has an even length, then the first half is a mirroring of the second half; e.g., \textit{abccba} is a palindrome, whereas \textit{abccab} is not. $(ii)$ if the string has an odd length, then disregard of the letter in the middle, the first half is a mirroring of the second half; e.g., \textit{abcvcba} is a palindrome, whereas \textit{abccvba} is not.\\
    
    \textbf{Hints}: 
    \begin{itemize}
        \item In your \textit{main()} when initializing \textit{mystr[]}, recall that a `\textbackslash0' will be included. To counter that, you'd probably need to create another string \textit{newstr[]} that only contains the actual letters.\smallskip
        
        \item In your \textit{Palindrome()} function, your \textit{base cases} should consider when the current string has either a length of 0 and 1.\smallskip
        
        \item In your \textit{Palindrome()} function, you want to match the first and last chars in your current string. What's the implication if that's not the case? And if that is the case (i.e., a match), how do we proceed with the next call to check the next pair (i.e., what is your \textit{recursive step})?\smallskip
        
        \item In your \textit{recursive step}, you will probably need to create another string \textit{newstr[]} for the next recursive call. Also what is the length \textit{len} you will want to pass in?\smallskip
        
        \item Your \textit{Palindrome()} function will take about 20-25 lines. If yours is way beyond those numbers, then you are probably doing something wrong.
    \end{itemize}
    
    
    
    
\end{enumerate}


\end{document}
